\section{Beam 1: organization of knowledge}

\subsection{Primary foundation}

The primary formal foundation is \citet{garicano2000}. Garicano studies production in which
workers encounter heterogeneous problems, knowledge acquisition is costly, and communication
allows a worker who does not know a solution to ask another worker.

Let
\[
Z\in\Q
\]
denote the problem encountered, with distribution \(F\) and density \(f\). A worker possesses
a knowledge set
\[
A\subseteq\Q,
\]
and solves the problem if \(Z\in A\).

In the autarkic case, Garicano orders problems from more to less frequent and considers a
knowledge interval \([0,Z_a]\). If the marginal cost of acquiring knowledge is \(c\), expected
net output is
\begin{equation}
E[y^a] = F(Z_a)-cZ_a,
\end{equation}
with interior first-order condition
\begin{equation}
f(Z_a)-c=0.
\end{equation}
Thus knowledge is acquired up to the point at which the marginal probability of solving an
additional problem equals its marginal acquisition cost.

\subsection{Communication and hierarchy}

Communication permits different workers to possess different knowledge sets. In Garicano's
model an organization assigns workers to classes, associates a knowledge set with each class,
and specifies whom a worker can ask when a solution is unknown. Asking for help consumes the
time of the worker being consulted. Denote this helping cost by \(h\).

The central trade-off is therefore between:
\begin{itemize}
\item acquiring knowledge locally, which is costly but avoids consultation; and
\item specializing knowledge and consulting other workers, which economizes on duplicated
knowledge acquisition but consumes communication and expert capacity.
\end{itemize}

Under the model assumptions, the resulting organization has a knowledge hierarchy:
production workers handle common problems and specialized problem solvers handle increasingly
exceptional problems. Information flows vertically and the organization exhibits management
by exception and a pyramidal structure.

This result is particularly relevant because the hierarchy is not imposed merely as an
organizational metaphor. It follows from the model's costs, problem frequencies and time
constraints.

\subsection{What Beam 1 establishes for the present skeleton}

Beam~1 supports the proposition that, for a given problem distribution and a given organization
of available knowledge, it can be inefficient for every operational unit to possess or use
all available expertise. Work can instead be organized so that common problems are handled
locally and exceptional problems are escalated.

In abstract notation, Beam~1 supplies a mapping of the form
\begin{equation}
a_t = \mathcal{R}(F_t,K_t,\kappa_t),
\label{eq:routing}
\end{equation}
where \(F_t\) describes problem demand, \(K_t\) represents the currently available organization
of knowledge, \(\kappa_t\) collects relevant costs and capacity constraints, and \(a_t\) is an
allocation/escalation organization.

Equation~\eqref{eq:routing} is an HLS-oriented abstraction, not an equation stated by
Garicano.

\subsection{What Beam 1 does not establish}

The base model explicitly assumes that the same problem recurs with probability zero and,
consequently, that communication does not itself generate learning. Garicano therefore does
not provide the temporal transition required for competence evolution.

Garicano also studies a knowledge set \(A\): whether solutions to classes of problems are
known. This should not be silently reinterpreted as a continuous measure of proficiency,
accuracy, speed or model quality.

Section IV of \citet{garicano2000} considers overlapping knowledge motivated by on-the-job
learning, but this does not constitute a dynamic law in which today's allocation endogenously
updates tomorrow's competence state. Beam~2 is required for that role.
